\chapter{Conclusão}
\label{ch:conc}

Através da realização deste projeto no âmbito da Unidade Curricular Sistemas Autónomos Inteligentes, desenvolveram-se técnicas e competências aprendidas durante o percurso universitário e conseguiu-se também entender de uma forma mais explícita e concreta como funciona o desenvolvimento de um projeto ligado à navegação autónoma de um robô.

Neste trabalho, criou-se um robô em CAD, numa plataforma nunca utilizada pelos elementos do grupo. Exportou-se esse robô para um ambiente de simulação, onde teve que passar por várias etapas até se obter um robô capaz de desempenhar uma navegação autónoma, através do uso do código que foi produzido e alterado de acordo com as características do robô e com o que se pretendia que ele fizesse. Por fim, ainda se adicionou mais um robô demonstrando que o código é adaptável e que se poderá futuramente ter uma série de robôs a navegar autonomamente e executando o mesmo trajeto. Uma outra sugestão consiste em testar o funcionamento destes robôs noutros cenários.

Pode-se assim afirmar que se enfrentaram algumas adversidades enquanto se desenvolvia o protótipo. Não obstante, tentou-se sempre encontrar soluções para resolver esses problemas e nos casos em que tal não foi possível, procurou-se sempre contorná-los. No final, obteve-se sucesso nos resultados alcançados visto que foram de encontro ao que era pretendido, tanto num robô como nos dois.

É de salientar ainda que todo o tempo utilizado por todos os membros do grupo foi gerido de forma a ser o mais proveitoso possível. 

Por último, são projetos como estes que colocam os alunos mais próximos do mundo do trabalho, uma vez que com a realização deste projeto, se sentiu que é fulcral recorrer-se às simulações para futuramente colocar-se em prática e obter-se um bom desempenho.
